\chapter{Kurzfassung}

Die genaue Vorhersage von Kryptowährungspreisen ist entscheidend um eine informierte Entscheidungsfindung für Händler und Investoren zu ermöglichen. Diese Arbeit präsentiert einen innovativen Ansatz, der fortschrittliche Techniken des maschinellen Lernens, umfassende Datenanalysen und eine starke technologische Basis kombiniert, um stündliche Schlusskurse von Kryptowährungen vorherzusagen. \\

Das System nutzt eine durchgehende Datenpipeline, die Marktinformationen in Echtzeit von der Kraken Kryptobörse aufzeichnet, vorverarbeitet und analysiert. Diese Daten werden in ein Open-, High-, Low-, Close-Preis (OHLC)-Format umgewandelt und um Volumen- sowie Handelszahlenmetriken ergänzt, um eine ganzheitliche Analyse möglichen zu machen. \\

Im Mittelpunkt dieses Systems wird ein Long Short-Term Memory (LSTM)-Algorithmus eingesetzt, der mit traditionellen Finanzindikatoren angereichert ist. Dieser wird trainiert und genau abgestimmt, um die stündlichen Schlusskurse vorherzusagen. Der Einsatz von TimescaleDB sorgt für eine effiziente Datenspeicherung und -abfrage, während MLflow die Überwachung der Modellmetriken und -architekturen ermöglicht und so den Grundstein für kontinuierliche Verbesserungen legt. Zusätzlich bietet MinIO eine skalierbare Lösung für die sichere Speicherung von Artefakten. \\

Durch eine Architektur, die ein Rust-basiertes Backend und einen Python-basierten Webserver kombiniert, wird eine hohe Leistung und die Echtzeitausführung des Modells gewährleistet. Eine benutzerfreundliche Oberfläche ermöglicht es Händlern und Investoren, intuitiv auf analytische Einblicke und Prognosen zuzugreifen. \\

Obwohl das System vielversprechende Ansätze aufweist, befindet es sich in einer Phase, in der kontinuierliche Verbesserungen möglich sind und bietet dadurch eine solide Grundlage für zukünftige Entwicklungen in der Erstellung von Modellen zur Vorhersage von Kryptowährungspreisen. Diese Forschungsarbeit leistet einen Beitrag zu diesem Feld und stellt ein ausgefeiltes Werkzeug-Set zur Verfügung, das Händlern, Investoren und Forschern helfen kann, ihren Entscheidungsprozess zu verbessern.